\chapter{Martingales, Strategies and Almost-Sure Convergence}
\label{sec:p15-martingales-and-convergence}

A martingale preserves its current conditional mean as information grows.
Submartingales and supermartingales replace that equality by opposite
conditional inequalities. The conditional calculus permits rigorous
transformations, stopping arguments and convergence proofs, provided their
measurability and moment hypotheses are retained.
The source for the constructions, theorems and exercises developed here
is Durrett's Section~4.2 \cite[pp.~217--224]{DurrettPTEAuthor2019}.

\section{Filtrations, processes and basic examples}
\label{sec:p15-martingale-definitions}

\begin{hbdefinition}{def:p15-filtration-martingale}{Information indexed by discrete time}
Fix \((\Omega,\mathcal F,P)\). A filtration is an increasing sequence
\((\mathcal F_n)_{n\geq0}\) of sub-sigma-algebras of \(\mathcal F\).
A real process \(X=(X_n)_{n\geq0}\) is adapted if \(X_n\) is
\(\mathcal F_n\)-measurable for every \(n\).
It is a martingale if each \(X_n\) is integrable and
\[
                         E(X_{n+1}\mid\mathcal F_n)=X_n
                                      \quad(n\geq0).
\]
It is a submartingale when the conditional expectation is at least \(X_n\),
and a supermartingale when it is at most \(X_n\).
These identities and inequalities hold almost everywhere at each time.
There are countably many times, so their exceptional sets may be combined.
The natural filtration of \(X\) is
\(\mathcal F_n^X=\sigma(X_0,\ldots,X_n)\); its initial value is included.
Write \(\Delta X_n=X_n-X_{n-1}\) for \(n\geq1\).
\index{filtration}\index{martingale}\index{submartingale}\index{supermartingale}
\end{hbdefinition}

\begin{hbtheorem}{thm:p15-martingale-examples}{proved=annex:p15.martingale-examples}{Linear, quadratic, product and exponential constructions}
Let \(\xi_i\) be i.i.d., \(E|\xi_i|<\infty\), \(E\xi_i=\mu\), and
\(S_n=c+\sum_{i=1}^n\xi_i\) for a real constant \(c\).
Use \(\mathcal F_n=\sigma(\xi_1,\ldots,\xi_n)\) and trivial \(\mathcal F_0\).
Then \(S_n-n\mu\) is a martingale; \(S_n\) is a submartingale for
\(\mu\geq0\) and a supermartingale for \(\mu\leq0\).
If \(\operatorname{Var}(\xi_i)=\sigma^2<\infty\), then
\[
                              (S_n-n\mu)^2-n\sigma^2
\]
is a martingale.

For independent nonnegative \(Y_i\) with \(EY_i=1\), let
\(\mathcal H_0=\{\varnothing,\Omega\}\) and
\(\mathcal H_n=\sigma(Y_1,\ldots,Y_n)\).
Then \(M_0=1,\ M_n=\prod_{i=1}^nY_i\) is a martingale for \((\mathcal H_n)\).
Identical laws are unnecessary for this product assertion.
In the i.i.d. increment model, if \(\theta\in\mathbb R\) and
\(\varphi(\theta)=E e^{\theta\xi_1}<\infty\), then
\(0<\varphi(\theta)<\infty\) and
\[
              \frac{e^{\theta(S_n-c)}}{\varphi(\theta)^n},
              \qquad
              \frac{e^{\theta S_n}}{\varphi(\theta)^n}
\]
are martingales for \((\mathcal F_n)\) with initial values 1 and
\(e^{\theta c}\), respectively.
The initial-value factor is retained explicitly.
\end{hbtheorem}

\begin{hbtheorem}{thm:p15-martingale-conditioning}{proved=annex:p15.martingale-conditioning}{Several time steps and less information}
If \(X\) is a supermartingale and \(m\leq n\), then
\[
                              E(X_n\mid\mathcal F_m)\leq X_m.
\]
For a submartingale the inequality reverses; for a martingale it is equality.
In particular their expectations have the corresponding monotonicity.
If \(\mathcal G_n\subseteq\mathcal F_n\) is another filtration to which
\(X\) is adapted, the same martingale, submartingale or supermartingale
property holds for \((\mathcal G_n)\). This includes the natural filtration.
\end{hbtheorem}

\begin{hbtheorem}{thm:p15-martingale-convex}{proved=annex:p15.martingale-convex}{Convex transformations and maxima}
If \(X\) is a martingale, \(\phi:\mathbb R\to\mathbb R\) is convex and
\(E|\phi(X_n)|<\infty\) at every time, then \(\phi(X_n)\) is a
submartingale. In particular \(|X_n|^p\) is a submartingale when
real \(p\geq1\) and its moments are finite.
For a submartingale \(X\), the same conclusion requires \(\phi\)
nondecreasing as well as convex.
Consequently \((X_n-a)^+\) is a submartingale for a submartingale \(X\),
and \(X_n\wedge a\) is a supermartingale for a supermartingale \(X\).
The maximum of two submartingales for the same filtration is a
submartingale; the minimum of two supermartingales is a supermartingale.
\end{hbtheorem}

\section{Predictable transforms and stopping}
\label{sec:p15-martingale-stopping}

\begin{hbdefinition}{def:p15-predictable-transform}{A strategy based on the previous information}
A real sequence \(H_n\), \(n\geq1\), is predictable if \(H_n\) is
\(\mathcal F_{n-1}\)-measurable. Its transform of \(X\) is
\[
                    (H\cdot X)_0=0,\qquad
                    (H\cdot X)_n=\sum_{i=1}^nH_i\Delta X_i.
\]
Here each \(H_i\) will be bounded by a deterministic finite constant
that may depend on \(i\). No uniform bound over all times is assumed.
Predictability excludes dependence on the increment presently being transformed.
\index{predictable process}\index{martingale transform}
\end{hbdefinition}

\begin{hbtheorem}{thm:p15-predictable-transform}{proved=annex:p15.predictable-transform}{The transform theorem}
For predictable bounded \(H_i\), the transform is adapted and integrable
at every finite time.
If \(X\) is a martingale, \(H\cdot X\) is a martingale, with either sign
of \(H_i\) allowed.
If \(H_i\geq0\), the transform preserves the submartingale or
supermartingale property as well.
\end{hbtheorem}

\begin{hbdefinition}{def:p15-stopping-time}{Stopping times and the information available when stopping}
A stopping time is a map \(T:\Omega\to\{0,1,\ldots\}\cup\{\infty\}\)
such that \(\{T\leq n\}\in\mathcal F_n\) for every \(n\).
Define
\[
 \mathcal F_T=\{A\in\mathcal F:
                       A\cap\{T\leq n\}\in\mathcal F_n\text{ for every }n\}.
\]
The stopped process is \(X^T_n=X_{T\wedge n}\).
It always uses a finite index; no value \(X_\infty\) is needed when \(T=\infty\).
Infima of empty sets of candidate stopping indices are defined to be \(\infty\).
\index{stopping time}\index{stopped process}
\end{hbdefinition}

\begin{hbtheorem}{thm:p15-stopping-calculus}{proved=annex:p15.stopping-calculus}{Stopping events, stopped information and stopped processes}
The stopping condition is equivalent to
\(\{T=n\}\in\mathcal F_n\) for each finite \(n\).
The collection \(\mathcal F_T\) is a sigma-algebra; if stopping times
satisfy \(S(\omega)\leq T(\omega)\) for every \(\omega\), then
\(\mathcal F_S\subseteq\mathcal F_T\).
For bounded \(T\), \(X_T\) is \(\mathcal F_T\)-measurable and integrable
when \(X\) is adapted and integrable at every finite time.
For any stopping time, stopping preserves all three martingale classes:
\[
                        X_{T\wedge n}
                            =X_0+\sum_{i=1}^n\mathbf1_{\{T\geq i\}}\Delta X_i.
\]
The indicators are predictable and bounded.
\end{hbtheorem}

\begin{hbtheorem}{thm:p15-bounded-optional-sampling}{proved=annex:p15.bounded-optional-sampling}{Optional sampling for bounded times}
If \(S\leq T\leq N\) are stopping times, with a deterministic
integer \(N\geq0\),
then for a martingale \(X\),
\[
                              E(X_T\mid\mathcal F_S)=X_S.
\]
For a supermartingale the left side is at most \(X_S\); for a
submartingale it is at least \(X_S\).
In particular \(EX_T=EX_0\) for bounded \(T\) and a martingale.
The conclusion at an unbounded time needs additional hypotheses.
\end{hbtheorem}

\section{Upcrossings and convergence}
\label{sec:p15-martingale-upcrossings}

\begin{hbdefinition}{def:p15-upcrossing-times}{The exact alternating stopping construction}
For real \(a<b\), set the auxiliary initial marker \(T_0=-1\), and define
for \(j\geq1\)
\[
 T_{2j-1}=\inf\{m>T_{2j-2}:X_m\leq a\},\qquad
 T_{2j}=\inf\{m>T_{2j-1}:X_m\geq b\},
\]
where candidate indices are nonnegative integers.
If a preceding time is infinite, every subsequent time is infinite.
Define
\[
 U_n(a,b)=\#\{j\geq1:T_{2j}\leq n\},\qquad
 U_\infty(a,b)=\lim_nU_n(a,b),
\]
and, for \(m\geq1\),
\[
 H_m=\mathbf1_{\{T_{2j-1}<m\leq T_{2j}\text{ for some }j\geq1\}}.
\]
The count includes a crossing begun at time zero. The marker \(T_0=-1\)
is only an initialization, not a stopping time in the nonnegative time set.
\index{upcrossing}
\end{hbdefinition}

\begin{hbtheorem}{thm:p15-upcrossing-inequality}{proved=annex:p15.upcrossing-inequality}{Path counting and the source upcrossing inequality}
The alternating times \(T_j\), \(j\geq1\), are stopping times, and
\(H_m\) is predictable with values zero or one.
For each deterministic finite path, \(U_n(a,b)\) is the maximum number
of chronologically disjoint low-to-high crossings completed by time \(n\).
If \(X\) is a submartingale, then
\[
                 (b-a)EU_n(a,b)
                      \leq E[(X_n-a)^+]-E[(X_0-a)^+].
\]
The proof uses the source's clipped process \(a+(X_n-a)^+\), so that
an unfinished crossing contributes a nonnegative amount.
\end{hbtheorem}

\begin{hbtheorem}{thm:p15-martingale-convergence}{proved=annex:p15.martingale-convergence}{Almost-sure convergence from a one-sided moment bound}
If \(X\) is a submartingale with \(C=\sup_nE(X_n^+)<\infty\),
then \(X_n\) converges almost surely to a finite integrable variable \(X_\infty\),
with
\[
              EX_\infty^+\leq C,\qquad
              EX_\infty^-\leq C-EX_0,\qquad
              E|X_\infty|\leq\liminf_nE|X_n|.
\]
The sign-reversed conclusion holds for a supermartingale with
\(\sup_nE(X_n^-)<\infty\).
A nonnegative supermartingale therefore converges almost surely to
an integrable nonnegative limit with \(EX_\infty\leq EX_0\) and
\(E(X_\infty\mid\mathcal F_n)\leq X_n\) for every \(n\).
In particular an \(L^1\)-bounded martingale has an integrable
almost-sure limit. These hypotheses alone do not imply \(L^1\) convergence.
\end{hbtheorem}

\begin{hbtheorem}{thm:p15-martingale-source-counterexamples}{proved=annex:p15.martingale-source-counterexamples}{The source's two distinct convergence failures}
The following two constructions have the stated properties.
\begin{enumerate}
\item Let \(S_n=1+\sum_{i=1}^n\xi_i\) for independent fair signs,
      \(T=\inf\{n:S_n=0\}\), and \(X_n=S_{T\wedge n}\).
      Then \(T<\infty\) almost surely, \(X\) is a nonnegative martingale,
      \(EX_n=1\) for every \(n\), and \(X_n\to0\) almost surely but
      \(E|X_n|=1\) does not tend to zero.
\item Let \(B_k\) be independent Bernoulli variables with \(P(B_k=1)=1/k\),
      and let \(\varepsilon_k\) be independent fair signs, independent
      of the whole \(B\)-sequence. Set \(X_0=0\) and
      \[
      X_k=
      \begin{cases}
       0,&B_k=0,\\
       \varepsilon_k,&B_k=1,\ X_{k-1}=0,\\
       kX_{k-1},&B_k=1,\ X_{k-1}\ne0 .
      \end{cases}
      \]
      For the filtration of these inputs, \(X\) is a martingale and
      \(P(X_k\ne0)=1/k\). It converges to zero in probability but
      has no almost-sure limit.
\end{enumerate}
The second construction supplies independent driving variables for
every transition, making the Borel--Cantelli step explicit.
\end{hbtheorem}

\section{Localization, products and switching}
\label{sec:p15-martingale-refinements}

\begin{hbtheorem}{thm:p15-localized-martingale-convergence}{proved=annex:p15.localized-martingale-convergence}{Pathwise upper bounds and integrable overshoots}
Let \(X\) be a submartingale with \(\sup_nX_n<\infty\) almost surely.
If
\[
                             E\left[\sup_{n\geq1}(\Delta X_n)^+\right]<\infty,
\]
then \(X_n\) converges almost surely to a finite limit.
The proof localizes at the first crossing above each positive integer
level and bounds the overshoot by the displayed integrable variable.
\end{hbtheorem}

\begin{hbtheorem}{thm:p15-nontrivial-product-limit}{proved=annex:p15.nontrivial-product-limit}{Nontrivial mean-one products and their logarithmic rate}
Let \(Y_i\) be i.i.d. nonnegative variables with \(EY_i=1\) and
\(P(Y_i=1)<1\). Then \(M_n=\prod_{i=1}^nY_i\) tends to zero almost surely.
For the natural logarithm, with \(\log0=-\infty\),
\[
                         \frac1n\log M_n\longrightarrow c
                                      \quad\text{almost surely},
 \qquad c=E\log Y_1\in[-\infty,0).
\]
The positive part of \(\log Y_1\) is integrable, so this extended
expectation is defined. The proof separately handles zero factors,
an integrable logarithm, and an infinite negative logarithmic mean.
\end{hbtheorem}

\begin{hbtheorem}{thm:p15-summable-product-compensation}{proved=annex:p15.summable-product-compensation}{Convergent products and a summable multiplicative error}
For deterministic \(y_n>-1\) with \(\sum_n|y_n|<\infty\),
the product \(\prod_n(1+y_n)\) converges to a finite strictly positive number.
Now let \(X_n,Y_n\) be nonnegative, integrable and adapted, with
\[
                E(X_{n+1}\mid\mathcal F_n)\leq(1+Y_n)X_n,
                \qquad \sum_{n\geq0}Y_n<\infty\quad\text{almost surely}.
\]
Set \(A_0=1,\ A_n=\prod_{j=0}^{n-1}(1+Y_j)\).
Then \(X_n/A_n\) is a nonnegative supermartingale, and \(X_n\)
converges almost surely to a finite limit.
No integrability of that final unnormalized limit is asserted from
these hypotheses alone.
\end{hbtheorem}

\begin{hbtheorem}{thm:p15-supermartingale-switching}{proved=annex:p15.supermartingale-switching}{The source switching principle}
Let \(X^1,X^2\) be supermartingales for the same filtration, and let
\(T\) be a stopping time such that \(X^1_T\geq X^2_T\) almost surely
on \(\{T<\infty\}\). Then both
\[
 Y_n=X^1_n\mathbf1_{\{T>n\}}+X^2_n\mathbf1_{\{T\leq n\}},
 \qquad
 Z_n=X^1_n\mathbf1_{\{T\geq n\}}+X^2_n\mathbf1_{\{T<n\}}
\]
are supermartingales. The two formulas switch at different discrete times,
and both are proved with those exact inequalities.
\end{hbtheorem}

\begin{hbtheorem}{thm:p15-dubins-upcrossing}{proved=annex:p15.dubins-upcrossing}{Dubins' bound}
For a nonnegative supermartingale \(X\), \(0<a<b\), and an integer \(k\geq1\),
\[
                 P(U_\infty(a,b)\geq k)
                     \leq\left(\frac ab\right)^k
                                  E\min\left(\frac{X_0}{a},1\right).
\]
The proof supplies the source's alternating scaled-process construction
and its stopped supermartingales using the switching principle.
The positive lower threshold and positive crossing count are part of
this bound.
\end{hbtheorem}

\section{Calculated strategy and hypothesis boundaries}
\label{sec:p15-martingale-boundaries}

\begin{hbtheorem}{thm:p15-martingale-boundary-calculations}{proved=annex:p15.martingale-boundary-calculations}{Predictability, unbounded stopping and missing moment control}
The following are fully specified mathematical models.
\begin{enumerate}
\item For a fair-sign random walk, the coefficients \(H_n=\xi_n\)
      give \(\sum_{i=1}^nH_i\xi_i=n\). They are bounded but not predictable,
      and the transform is not a martingale.
\item For independent signs with \(P(\xi_n=1)=p>0\), let \(H_1=1\),
      and double the previous stake after a loss, resetting it to one
      after a success. At the first success \(T\), the transform gain is one.
      Before that stopping time, at horizon \(n\), it is
      \[
        G_n=\begin{cases}1,&T\leq n,\\1-2^n,&T>n.\end{cases}
        \qquad EG_n=1-[2(1-p)]^n.
      \]
      At \(p=1/2\), every finite-horizon expectation is zero while
      \(G_n\to1\) almost surely and \(E|G_n-1|=1\).
\item Independent variables \(\xi_n\) with value \(n^2-1\) with
      probability \(1/n^2\) and value \(-1\) otherwise have mean zero.
      Their sums form a martingale that tends to \(-\infty\)
      almost surely; in fact \(n^{-1}\sum_{i=1}^n\xi_i\to-1\).
\item The deterministic sequence \(X_n=-1/(n+1)\) is a submartingale,
      whereas \(X_n^2\) is a supermartingale. Convexity alone does
      not preserve the submartingale property.
\end{enumerate}
For \(0<p<1\), the first-success strategy has a finite almost-sure
stopping time but no deterministic time bound. The fair-case limit
does not satisfy the hypotheses of bounded optional sampling.
\end{hbtheorem}

\begin{proofboundary}{def:p15-martingale-next}{Stronger modes of convergence and unbounded sampling}
Almost-sure convergence holds under the stated one-sided moment bounds.
The uniform-integrability criterion and the forward/backward conditional
limits are supplied in
\cref{thm:p15-vitali-convergence,thm:p15-closed-martingale,thm:p15-levy-upward,thm:p15-levy-backward}.
The full maximal estimates and \(L^p\) convergence results are
\cref{thm:p15-doob-lp-maximal,thm:p15-martingale-lp-convergence},
with the \(L^1\) boundary in \cref{thm:p15-fair-walk-maximal-tail}.
For unbounded almost-surely finite stopping times, the square-root
variation criterion and the conditional sampling conclusion are
\cref{thm:p15-square-root-stopping}; their moment hypotheses are
essential, as the fair-walk hitting time in that theorem shows.
The full Section~4.8 unbounded and infinite-time sampling results are
\cref{thm:p15-os-ui-stopped,thm:p15-os-conditional-sampling,thm:p15-os-nonnegative}, with both Wald identities and their precise
moment conditions in \cref{thm:p15-os-wald-first,thm:p15-os-wald-second}.
The source's separate potential-theoretic explanation of the terminology
is not a premise of these discrete-time results.
\end{proofboundary}
